% section2_bayesian_surprise.tex

\section{Psychological Exposure and Bayesian Surprise}
\label{sec:surprise}

The benchmark environment in Section~\ref{sec:model} treats the ambitious signal $s_A$ as a costly signal of ability. This section adds the psychological object that distinguishes the model from standard costly signaling. The ambitious signal does not only reveal information about the sender's type. It also exposes the sender's self-model. A sender who chooses $s_A$ represents herself as someone whose private assessment of her own capacity is accurate enough to justify taking the ambitious action. The social cost of later reinterpretation is therefore not only a material or reputational loss. It is also an epistemic loss: observers revise their beliefs about whether the sender correctly understood her own type.

The section makes three clarifications. First, the exposure index $\Phi_\theta$ is a true-type-indexed primitive object, derived from primitive self-model-accuracy levels. It is not an equilibrium object generated by a previous separating history. Second, the relevant channel is adverse repricing, not truth verification. A high-precision audit would tend to confirm genuinely accurate high types; that is not the technology studied here. Third, the D1 argument later uses $\Phi_\theta$ as a cost attached to the sender's true type, not as a cost attached to the receiver's off-path attribution. This keeps the exposure technology separate from the belief refinement that interprets deviations.

\subsection{Belief-dependent exposure}

Let $j\in\{L,H\}$ denote the current judgment-intensity state, with associated scrutiny multiplier $\mu_j>0$. The low state $L$ represents ordinary evaluation and the high state $H$ represents intensified evaluation. The state variable $j$ is not a political primitive. It is an exogenous state of the social evaluation technology.

After $s_A$, observers form beliefs not only about the sender's ability $\theta$, but also about the accuracy of the sender's self-model. Write
\[
m\in\{0,1\}
\]
for the event that the sender's self-model is accurate, where $m=1$ denotes accurate self-assessment. The sender cares not only about what observers infer about $\theta$, but also about what observers infer about the sender's own assessment of $\theta$.

For a finite hierarchy of beliefs $(\mu^2,\ldots,\mu^k)$, let
\[
L(\mu^2,\ldots,\mu^k;j,x)
\]
denote the epistemic penalty in judgment state $j$ and exposure state $x\in\{0,1\}$. The exposure state records whether the sender previously chose the ambitious signal. Choosing $s_A$ sets $x=1$ next period; choosing $s_P$ clears exposure next period.

The total one-period psychological payoff can be written as
\[
u_\theta(s,j,x;\mu^2,\ldots,\mu^k)
=
v_\theta(s)-c_\theta(s)-\mu_j L(\mu^2,\ldots,\mu^k;j,x),
\]
where $v_\theta(s)$ is the ordinary reputational or material payoff from the signal, $c_\theta(s)$ is the intrinsic cost of the signal, and the final term is the judgment-weighted epistemic penalty.

The next assumption is purely regularity. It rules out discontinuities inside a fixed judgment state. The discontinuities studied below will come from changes in the transition probability into high scrutiny, not from discontinuous psychological utility within a regime.

\begin{assumption}[Regime-conditional continuity]
\label{ass:conditional-continuity}
For each fixed $j\in\{L,H\}$ and $x\in\{0,1\}$, the function
\[
L(\mu^2,\ldots,\mu^k;j,x)
\]
is continuous in the relevant finite hierarchy of higher-order beliefs. The action set, type space, and exposure state space are finite.
\end{assumption}

\begin{lemma}[Regime-conditional continuity]
\label{lem:regime-continuity}
Under Assumption~\ref{ass:conditional-continuity}, for each fixed judgment-intensity state $j$, the induced psychological payoff is continuous in the relevant higher-order belief hierarchy. Hence the regime-conditional psychological game satisfies the standard continuity requirement for belief-dependent payoffs. On any open region where the relevant incentive constraints are strict, the associated pure-strategy equilibrium selection is locally constant.
\end{lemma}

\begin{proof}
Fix $j$ and $x$. The ordinary payoff terms $v_\theta(s)$ and $c_\theta(s)$ are defined on a finite action and type space. The multiplier $\mu_j$ is constant within the state $j$. By Assumption~\ref{ass:conditional-continuity}, the epistemic penalty $L(\mu^2,\ldots,\mu^k;j,x)$ is continuous in the relevant belief hierarchy. Multiplication by the fixed scalar $\mu_j$ and addition of finite ordinary payoff terms preserve continuity. Thus, the regime-conditional psychological payoff is continuous.

Since the type and action spaces are finite, strict incentive inequalities are preserved under sufficiently small perturbations of beliefs and payoff parameters within a fixed state. Therefore, wherever the relevant incentive constraints are strict, the same pure-strategy profile remains selected in a neighborhood of the primitive. Any discontinuous change in equilibrium behavior must therefore come from a change in the relevant state, transition probability, or binding incentive constraint, rather than from a discontinuity of belief-dependent utility within the fixed state.
\end{proof}

\subsection{Primitive self-accuracy levels}

The exposure index is built from a level channel, not a precision channel. Let
\[
p_\theta:=\Pr(m=1\mid \theta)
\]
denote the primitive probability that a type-$\theta$ sender has an accurate self-model. The key ordering is
\[
1-\delta>p_H>p_L>\bar p>\delta>0.
\]
Higher types are genuinely more self-accurate in the primitive environment. The lower benchmark $\bar p$ is the assessment induced by adverse repricing under high scrutiny.

This assumption should not be read as signal precision. Precision is a fidelity concept: a more precise audit reveals the true state more accurately. If high types are genuinely more accurate, a fidelity-improving audit tends to confirm them, not punish them. The model instead studies adverse social repricing. In a high-scrutiny state, an exposed self-presentation can be reinterpreted downward with a positive probability that is not driven by the sender's true merit. The economic loss is increasing in the amount of self-model reputation-at-risk carried by the true type.

Let $\alpha\in(0,1]$ be the probability, conditional on high scrutiny and exposure, that the ambitious self-presentation is adversely repriced. This probability is taken to be independent of type in the baseline. It can be normalized into $\mu_H$, but keeping it visible makes the interpretation clear: high scrutiny is not a deterministic truth audit. It is a state in which exposed self-presentations face a type-independent risk of adverse reinterpretation.

The epistemic exposure index is
\[
\Phi_\theta
=
\alpha\,
D_{\mathrm{KL}}
\left(
\operatorname{Bern}(p_\theta)
\,\middle\|\,
\operatorname{Bern}(\bar p)
\right).
\]
The direction of the divergence reflects the economic object. The loss is not the ordinary probability of failure. It is the collapse from a carried self-accuracy level to a lower adverse benchmark. Judgment intensity prices this exposure:
\[
C_E(\theta,j)=\mu_j\Phi_\theta.
\]

Thus $\Phi_\theta$ is neither history-induced reputation capital nor signal precision. It is the true-type-indexed loss from adverse repricing of self-model exposure. The receiver's later off-path attribution may affect the reputational return to a deviation, but it does not redefine $\Phi_\theta$.

\subsection{Bayesian-surprise complementarity}

\begin{lemma}[Bayesian-surprise foundation for scrutiny-type complementarity]
\label{lem:bayesian-surprise-complementarity}
Suppose
\[
1-\delta>p_H>p_L>\bar p>\delta>0.
\]
Define
\[
\Phi_\theta
=
\alpha\,
D_{\mathrm{KL}}
\left(
\operatorname{Bern}(p_\theta)
\,\middle\|\,
\operatorname{Bern}(\bar p)
\right)
\]
with $\alpha>0$, and let
\[
C_E(\theta,j)=\mu_j\Phi_\theta.
\]
Then the epistemic cost has increasing differences in type and judgment intensity. For every $\mu'>\mu$,
\[
C_E(\theta_H,\mu')-C_E(\theta_H,\mu)
>
C_E(\theta_L,\mu')-C_E(\theta_L,\mu).
\]
Equivalently, the marginal cost of higher judgment intensity is larger for the high-self-accuracy type.
\end{lemma}

\begin{proof}
For Bernoulli beliefs,
\[
D_{\mathrm{KL}}
\left(
\operatorname{Bern}(p)
\,\middle\|\,
\operatorname{Bern}(\bar p)
\right)
=
p\log\frac{p}{\bar p}
+
(1-p)\log\frac{1-p}{1-\bar p}.
\]
Differentiating with respect to $p$ gives
\[
\frac{\partial}{\partial p}
D_{\mathrm{KL}}
\left(
\operatorname{Bern}(p)
\,\middle\|\,
\operatorname{Bern}(\bar p)
\right)
=
\log\frac{p(1-\bar p)}{\bar p(1-p)}.
\]
This derivative is positive whenever $p>\bar p$. Since $p_H>p_L>\bar p$, it follows that
\[
\Phi_H>\Phi_L.
\]
Therefore, for every $\mu'>\mu$,
\[
\begin{aligned}
&\big[C_E(\theta_H,\mu')-C_E(\theta_H,\mu)\big]
-\big[C_E(\theta_L,\mu')-C_E(\theta_L,\mu)\big]\
&\qquad
=(\mu'-\mu)(\Phi_H-\Phi_L)>0.
\end{aligned}
\]
Thus higher judgment intensity raises the epistemic cost of the high-self-accuracy type more than that of the low-self-accuracy type.
\end{proof}

Lemma~\ref{lem:bayesian-surprise-complementarity} is the first reversal relative to standard signaling theory. In a Spence environment, high types are more willing to bear the ambitious signal because they face a lower ordinary cost of sending it. Here, high types also carry more self-model reputation-at-risk. When high scrutiny adversely reprices exposed self-presentation, the same type characteristic that supports ambitious action under low scrutiny becomes the source of larger exposure.

\begin{remark}[Repricing, not precision]
The sign of Lemma~\ref{lem:bayesian-surprise-complementarity} depends on the level channel. Replacing $p_\theta$ with signal precision would change the model. A higher-precision audit is a fidelity technology and tends to confirm an accurate high type. It therefore does not generate top-down withdrawal. The primitive needed here is a type-indexed self-accuracy level combined with an adverse repricing state, not a variance or precision parameter.
\end{remark}

\subsection{From self-accuracy to crash risk}

The preceding lemma converts primitive self-accuracy into exposure. A high type has more to gain from separation, but also more to lose from a future state that adversely reinterprets her ambitious signal as evidence about self-model risk.

Let $B_\theta$ denote the current net separation rent from choosing $s_A$ rather than $s_P$ in the low-judgment state, before future high-scrutiny exposure is priced. Let
\[
q=\Pr(j_{t+1}=H\mid j_t=L)
\]
be the transition probability into high scrutiny. If choosing $s_A$ creates the one-period shadow, the low-state payoff gain from choosing $s_A$ can be written as
\[
D_\theta^L(q)=B_\theta-\beta q\mu_H\Phi_\theta.
\]
The associated threshold is
\[
q_\theta^*=\frac{B_\theta}{\beta\mu_H\Phi_\theta},
\]
whenever $B_\theta>0$.

Thus top-down withdrawal is governed not by exposure alone, but by the ratio between exposure and current rent:
\[
\frac{\Phi_\theta}{B_\theta}.
\]
The high type exits first precisely when
\[
\frac{\Phi_H}{B_H}>\frac{\Phi_L}{B_L}.
\]
In that case,
\[
q_H^*<q_L^*.
\]
As the probability of future high scrutiny increases, the high type's incentive constraint binds first.

\begin{corollary}[Top-down threshold ordering]
\label{cor:top-down-threshold-ordering}
Suppose $B_H>0$, $B_L>0$, and
\[
\frac{\Phi_H}{B_H}>\frac{\Phi_L}{B_L}.
\]
Then
\[
q_H^*<q_L^*.
\]
Hence an increase in the transition probability into high scrutiny first destroys the ambitious-signal incentive of the type with the larger exposure-to-rent ratio.
\end{corollary}

\begin{proof}
The threshold is
\[
q_\theta^*=\frac{B_\theta}{\beta\mu_H\Phi_\theta}.
\]
Since $\beta\mu_H>0$,
\[
q_H^*<q_L^*
\]
iff and only if
\[
\frac{B_H}{\Phi_H}<\frac{B_L}{\Phi_L},
\]
which is equivalent to
\[
\frac{\Phi_H}{B_H}>\frac{\Phi_L}{B_L}.
\]
\end{proof}

The corollary gives the transition from the psychological exposure mechanism to the dynamic pooling result. The high type is not less capable of sending the ambitious signal. She is more exposed because her true type carries more self-model reputation-at-risk under adverse repricing.
